\documentclass[11pt]{article}
\usepackage[margin=1in]{geometry}
\usepackage{amsmath,amssymb,mathtools,amsthm}
\usepackage{booktabs,longtable}
\usepackage{enumitem}
\usepackage[hidelinks]{hyperref}

\newcommand{\dd}{\mathrm{d}}
\newcommand{\R}{\mathbb{R}}
\newcommand{\TT}{\mathrm{TT}}
\newcommand{\calX}{\mathcal{X}}
\newcommand{\calU}{\mathcal{U}}
\newcommand{\calD}{\mathcal{D}}
\newcommand{\calB}{\mathcal{B}}
\newcommand{\dotp}{\partial_i}

\newtheorem{proposition}{Proposition}
\newtheorem{corollary}{Corollary}
\theoremstyle{definition}
\newtheorem{definition}{Definition}
\newtheorem{assumption}{Assumption}
\theoremstyle{remark}
\newtheorem{remark}{Remark}

\title{A Human-Readable Derivation of Fixed-Branch Zhao--Cui Tensor-Train Filtering}
\author{BayesFilter P13 Human-Readable Note}
\date{2026-05-31}

\begin{document}
\maketitle

\begin{abstract}
This note explains, from first principles, the part of the Zhao--Cui
tensor-train sequential filtering construction that BayesFilter can use as a
candidate high-dimensional nonlinear filtering method.  The presentation is
intended for a reader who knows probability and numerical linear algebra, but
has not studied tensor-train filtering or the companion code.  We start with
ordinary filtering, work through a scalar nonlinear example, introduce tensor
trains only after the filtering object is clear, and then state two
propositions.  The first proposition proves that the fixed-branch squared-TT
recursion still produces normalized approximate filters.  The second proves
that its analytical gradient differentiates the same approximate likelihood
scalar computed by the fixed branch.  The result is a mathematical
construction, not an accuracy theorem: it says exactly what is computed and how
to differentiate it when the numerical branch is held fixed.
\end{abstract}

\section{What Problem Are We Solving?}

A state-space model describes a hidden state \(x_t\), a possible static
parameter \(\vartheta\), and an observed datum \(y_t\).  At time \(t\), the
filtering problem is to update the density of \((x_t,\vartheta)\) after seeing
\(y_1,\ldots,y_t\).  If the previous filter is known, the update is just
Bayes' rule, but in high dimension the integrals are the hard part.

Let
\[
  p_0(x_0,\vartheta),\qquad
  f_\alpha(x_t\mid x_{t-1},\vartheta),\qquad
  g_\alpha(y_t\mid x_t,\vartheta)
\]
be the prior, transition density, and observation density.  The symbol
\(\alpha\) denotes an external model parameter whose likelihood derivative may
be needed later.  The symbol \(\vartheta\) is a random coordinate carried by
the filter.

Suppose the previous filtering density is
\[
  p_{t-1}(x_{t-1},\vartheta\mid y_{1:t-1}).
\]
Before seeing \(y_t\), the model predicts
\[
  p(x_t,\vartheta\mid y_{1:t-1})
  =
  \int
  f_\alpha(x_t\mid x_{t-1},\vartheta)
  p_{t-1}(x_{t-1},\vartheta\mid y_{1:t-1})
  \,\dd x_{t-1}.
\]
After observing \(y_t\), Bayes' rule multiplies this prediction by
\(g_\alpha(y_t\mid x_t,\vartheta)\) and divides by the normalizing constant.

It is useful to keep the old state inside the calculation.  Define
\[
  u_t=(x_t,\vartheta,x_{t-1})
\]
and
\[
  q_t(u_t;\alpha)
  =
  p_{t-1}(x_{t-1},\vartheta\mid y_{1:t-1})
  f_\alpha(x_t\mid x_{t-1},\vartheta)
  g_\alpha(y_t\mid x_t,\vartheta).
  \label{eq:p13-exact-q}
\]
This is an unnormalized joint density over current state, parameter, and old
state.  Its integral is
\[
  Z_t(\alpha)
  =
  \int q_t(x_t,\vartheta,x_{t-1};\alpha)\,
  \dd x_t\,\dd\vartheta\,\dd x_{t-1}.
  \label{eq:p13-exact-evidence-increment}
\]
To see why this is the one-step evidence increment, integrate the joint density
in the order \(x_t,\vartheta,x_{t-1}\):
\[
\begin{aligned}
  Z_t(\alpha)
  &=
  \int
  g_\alpha(y_t\mid x_t,\vartheta)
  f_\alpha(x_t\mid x_{t-1},\vartheta)
  p_{t-1}(x_{t-1},\vartheta\mid y_{1:t-1})
  \,\dd x_t\,\dd\vartheta\,\dd x_{t-1}  \\
  &=
  \int
  g_\alpha(y_t\mid x_t,\vartheta)
  p(x_t,\vartheta\mid y_{1:t-1})
  \,\dd x_t\,\dd\vartheta
  =
  p_\alpha(y_t\mid y_{1:t-1}).
\end{aligned}
\]
The updated filter is the marginal of the normalized joint density:
\[
  p_t(x_t,\vartheta\mid y_{1:t})
  =
  \int
  \frac{q_t(x_t,\vartheta,x_{t-1};\alpha)}{Z_t(\alpha)}
  \,\dd x_{t-1}.
  \label{eq:p13-exact-filter-marginal}
\]
Thus a filtering algorithm needs three operations: approximate \(q_t\), compute
its integral, and marginalize out the old state.  Zhao and Cui's TT method is
one high-dimensional way to do exactly those three operations.

\section{A Scalar Example Before Tensor Trains}

Consider the one-dimensional nonlinear model
\[
  x_t=\rho x_{t-1}+\eta_t,\qquad
  y_t=x_t^2+\epsilon_t,
\]
where \(\eta_t\sim N(0,\sigma_\eta^2)\) and
\(\epsilon_t\sim N(0,\sigma_\epsilon^2)\).  Suppose the previous approximate
filter has density \(\widehat p_{t-1}(x_{t-1})\).  The unnormalized two-variable
object for one filtering step is
\[
  q_t(x_t,x_{t-1})
  =
  \widehat p_{t-1}(x_{t-1})
  \frac{1}{\sigma_\eta}
  \varphi\!\left(\frac{x_t-\rho x_{t-1}}{\sigma_\eta}\right)
  \frac{1}{\sigma_\epsilon}
  \varphi\!\left(\frac{y_t-x_t^2}{\sigma_\epsilon}\right),
  \label{eq:p13-scalar-q}
\]
where \(\varphi\) is the standard normal density.

The observation \(y_t=x_t^2+\epsilon_t\) makes this a useful toy example.  If
\(y_t\) is positive and the noise is small, the likelihood is large near
\(x_t\approx \sqrt{y_t}\) and \(x_t\approx-\sqrt{y_t}\).  The posterior can be
two-peaked even when the transition is Gaussian.  A single Gaussian projection
may miss that shape.  A square-root density approximation tries to represent
the non-Gaussian object more directly.

Let \(\phi_t(x_t,x_{t-1})\) be an approximation to \(\sqrt{q_t(x_t,x_{t-1})}\).
The squared approximation is
\[
  \widehat q_t(x_t,x_{t-1})
  =
  \phi_t(x_t,x_{t-1})^2+\tau_t\lambda_t(x_t,x_{t-1}),
  \label{eq:p13-scalar-qhat}
\]
where \(\tau_t\ge0\), \(\lambda_t\ge0\), and
\[
  \iint \lambda_t(x_t,x_{t-1})\,\dd x_t\,\dd x_{t-1}=1.
\]
The added term is a small defensive density.  It prevents the approximation
from assigning exactly zero mass in places where the numerical square-root fit
is unreliable.

The approximate evidence increment is
\[
  \widehat Z_t
  =
  \iint
  \widehat q_t(x_t,x_{t-1})\,\dd x_t\,\dd x_{t-1}
  =
  \iint \phi_t(x_t,x_{t-1})^2\,\dd x_t\,\dd x_{t-1}
  +\tau_t.
  \label{eq:p13-scalar-zhat}
\]
The approximate joint density after observing \(y_t\) is
\[
  \widehat\pi_t(x_t,x_{t-1})
  =
  \frac{\widehat q_t(x_t,x_{t-1})}{\widehat Z_t}.
\]
The next filter is its marginal over the old state:
\[
  \widehat p_t(x_t)
  =
  \int
  \widehat\pi_t(x_t,x_{t-1})\,\dd x_{t-1}.
  \label{eq:p13-scalar-filter}
\]
For a completely explicit defensive term, choose two normalized one-dimensional
densities \(\lambda_x\) and \(\lambda_{x^-}\), and set
\[
  \lambda_t(x_t,x_{t-1})
  =
  \lambda_x(x_t)\lambda_{x^-}(x_{t-1}).
\]
Then
\[
\begin{aligned}
  \widehat p_t(x_t)
  &=
  \frac{1}{\widehat Z_t}
  \int \phi_t(x_t,x_{t-1})^2\,\dd x_{t-1}
  +
  \frac{\tau_t}{\widehat Z_t}
  \lambda_x(x_t)
  \int \lambda_{x^-}(x_{t-1})\,\dd x_{t-1}\\
  &=
  \frac{1}{\widehat Z_t}
  \int \phi_t(x_t,x_{t-1})^2\,\dd x_{t-1}
  +
  \frac{\tau_t}{\widehat Z_t}\lambda_x(x_t).
\end{aligned}
  \label{eq:p13-scalar-filter-decomposition}
\]
This line closes the scalar recursion: fit \(\phi_t\), add the defensive
product density, normalize by \(\widehat Z_t\), and integrate out the old
state.
If \(\phi_t^2\) is close to \(q_t\), this filter should be a useful
approximation.  Even if it is not close, the algebra above still gives a
nonnegative normalized density.  That distinction matters: normalization is a
mathematical property of the construction; accuracy is a separate numerical
question.

\section{From The Scalar Example To Tensor Trains}

In the scalar example, \(\phi_t\) is a function of two variables.  In a high
dimensional state-space model, \(u_t=(x_t,\vartheta,x_{t-1})\) may have tens or
hundreds of coordinates.  A full tensor-product approximation over all
coordinates is usually impossible because the number of coefficients grows
exponentially.

A tensor train writes a high-dimensional function as a chain of smaller
matrix-valued functions.  Let \(u=(u_1,\ldots,u_D)\).  A functional TT
representation is
\[
  \phi(u)
  =
  G_1(u_1)G_2(u_2)\cdots G_D(u_D),
  \label{eq:p13-tt-eval}
\]
where
\[
  G_k(u_k)\in\R^{r_{k-1}\times r_k},
  \qquad r_0=r_D=1.
\]
The integers \(r_k\) are TT ranks.  Each \(G_k\) is a small matrix as a
function of one coordinate.  Multiplying the matrices yields a scalar.

For computation, each core is expanded in a one-dimensional basis:
\[
  G_k(u_k)
  =
  \sum_{a=1}^{n_k} C_{k,a} b_{k,a}(u_k),
  \qquad
  C_{k,a}\in\R^{r_{k-1}\times r_k}.
  \label{eq:p13-core-expansion}
\]
If the ranks \(r_k\) stay moderate, the number of stored coefficients scales
like \(\sum_k r_{k-1}n_k r_k\), rather than \(\prod_k n_k\).  That is the
basic high-dimensional promise: the representation is efficient when the
target density has low-rank structure in a suitable coordinate order and basis.

The squared representation uses this TT for a square-root density.  Once
\(\phi\) is represented by \eqref{eq:p13-tt-eval}, the approximation
\(\phi^2+\tau\lambda\) is automatically nonnegative.  The integral of
\(\phi^2\) can also be computed by contractions.  If
\[
  M_k[a,b]=\int b_{k,a}(u_k)b_{k,b}(u_k)\,\dd\nu_k(u_k),
\]
then define
\[
  R_D=1,\qquad
  R_{k-1}
  =
  \sum_{a,b=1}^{n_k}
  M_k[a,b]\,C_{k,a}R_kC_{k,b}^{\top},
  \qquad k=D,\ldots,1.
  \label{eq:p13-mass-contraction}
\]
The scalar \(R_0\) is exactly \(\int\phi(u)^2\,\dd\nu(u)\) for the declared
cores and mass matrices.  The approximation is in the fitted TT cores; the
contraction itself is just algebra.

\section{Zhao--Cui Algorithm In Human Language}

The Zhao--Cui method applies the square-root TT idea sequentially.  At each
time \(t\), it starts with the previous approximate filter, forms the joint
object \(q_t\) from \eqref{eq:p13-exact-q}, fits a TT approximation to a
square root of that object, squares the TT to obtain a nonnegative density,
normalizes it, and marginalizes out \(x_{t-1}\).  Written in BayesFilter
notation, the step is:
\[
  q_t(x_t,\vartheta,x_{t-1})
  =
  \widehat p_{t-1}(x_{t-1},\vartheta)
  f_\alpha(x_t\mid x_{t-1},\vartheta)
  g_\alpha(y_t\mid x_t,\vartheta).
\]
The fitted square-root TT is
\[
  \sqrt{q_t(u_t)}\approx
  \phi_t(u_t)
  =
  G_{t,1}(u_{t,1})\cdots G_{t,D}(u_{t,D}).
\]
The nonnegative approximation is
\[
  \widehat q_t(u_t)
  =
  \phi_t(u_t)^2+\tau_t\lambda_t(u_t),
  \qquad
  \int\lambda_t(u_t)\,\dd u_t=1.
\]
The approximate normalizer is
\[
  \widehat Z_t=\int\widehat q_t(u_t)\,\dd u_t,
\]
and the next filter is
\[
  \widehat p_t(x_t,\vartheta)
  =
  \int
  \frac{\widehat q_t(x_t,\vartheta,x_{t-1})}{\widehat Z_t}
  \,\dd x_{t-1}.
\]

This is the filtering connection in one sentence: the TT is not replacing the
state-space model; it is a deterministic representation used to approximate
the joint density whose normalization and marginalization are the filtering
update.

Conditional Knothe--Rosenblatt maps enter after the squared-TT density has
been built.  Because the squared TT can be integrated coordinate by coordinate,
one can form marginal and conditional distributions and then sample by
inverting conditional cumulative distribution functions.  That transport
machinery is important for sampling, smoothing, and sequential state-parameter
learning, but the two propositions below only need the density,
normalization, and marginalization pieces.  Zhao and Cui present the sequential
TT construction, squared-TT density, marginalization, and conditional transport
machinery in their state-space formulation; Cui and Dolgov provide the
underlying squared inverse Rosenblatt transport construction
\cite{zhao-cui-2024,cui-dolgov-2022}.

\section{Why The Published Adaptive Algorithm Is Not Yet A Gradient Algorithm}

For filtering alone, adaptation is natural.  A good TT implementation may
change ranks, interpolation points, active subspaces, local approximation
points, or truncation decisions as the data and parameter values change.  These
choices are part of how the method stays efficient.

A likelihood gradient has a stricter requirement.  It must be the derivative
of one scalar formula.  If a small change in \(\alpha\) changes the selected
rank, pivot set, sample set, truncation, or linear algebra branch, then the
formula being differentiated has changed.  The derivative of yesterday's
branch is still useful as a local sensitivity, but it is not an ordinary global
derivative of the adaptive algorithm.

This does not make the adaptive method wrong.  It only separates two uses:
\begin{itemize}[leftmargin=*]
\item For filtering and sampling, adaptive TT construction may be the right
  numerical tool.
\item For a same-scalar analytical gradient, we must either freeze the branch
  or design a smooth adaptation rule and differentiate that smooth rule.
\end{itemize}
The rest of this note follows the first option.  We freeze the numerical
branch so the value and derivative refer to the same approximate likelihood.

\section{The Fixed-Branch Algorithm}

A fixed branch is a fully specified deterministic version of one TT filtering
path.  It freezes every discrete or branch-changing choice that would otherwise
make the scalar formula change while differentiating.  The fixed branch does
not claim to be more accurate than the adaptive algorithm.  Its purpose is to
make the approximate likelihood and its analytical gradient unambiguous.

At each time \(t\), choose the following objects before differentiating:
\begin{enumerate}[leftmargin=*]
\item the coordinate order for \(u_t=(x_t,\vartheta,x_{t-1})\);
\item the reference domain or coordinate transform for each coordinate;
\item the one-dimensional bases \(b_{k,a}\) and mass matrices \(M_k\);
\item the TT ranks \(r_{t,k}\);
\item the interpolation, cross, or least-squares points used to fit the cores;
\item the local solver schedule, including sweep counts and stopping rules
  that are treated as fixed;
\item the defensive density \(\lambda_t\) and defensive mass \(\tau_t\);
\item any stabilizing constant \(c_t\) used to shift the fitted potential;
\item the signs and active ranks for QR, SVD, Cholesky, or related
  factorizations.
\end{enumerate}
After these choices are fixed, each step is an ordinary differentiable
calculation, provided the model densities and linear solves are differentiable
on that branch.

\paragraph{Forward step.}
Given \(\widehat p_{t-1}^{\TT}\), form
\[
  q_t(u_t;\alpha)
  =
  \widehat p_{t-1}^{\TT}(x_{t-1},\vartheta;\alpha)
  f_\alpha(x_t\mid x_{t-1},\vartheta)
  g_\alpha(y_t\mid x_t,\vartheta).
\]
If a stabilizing shift is used, write the square-root target evaluations as
\[
  s_t(u;\alpha)
  =
  \exp\!\left(-\frac12[V_{t,0}(u;\alpha)-c_t(\alpha)]\right),
\]
where \(V_{t,0}=-\log q_t\).  Equivalently,
\(s_t=\exp(c_t/2)\sqrt{q_t}\).  Fit TT cores
\[
  \phi_t(u;\alpha)=G_{t,1}(u_1;\alpha)\cdots G_{t,D}(u_D;\alpha)
\]
to these fixed target evaluations using the chosen interpolation or
least-squares equations.

Compute the square mass \(R_{t,0}\) by \eqref{eq:p13-mass-contraction}.  Then
\[
  \widehat z_t(\alpha)=R_{t,0}(\alpha)+\tau_t(\alpha)
\]
for fixed normalized \(\lambda_t\).  The original-scale likelihood increment
represented by the shifted fit is
\[
  \widehat L_t^{\TT}(\alpha)
  =
  \exp[-c_t(\alpha)]\,\widehat z_t(\alpha).
\]
The declared log-likelihood contribution is therefore
\[
  \widehat\ell_t^{\TT}(\alpha)
  =
  \log\widehat z_t(\alpha)-c_t(\alpha).
  \label{eq:p13-declared-step-scalar}
\]
Normalize and marginalize:
\[
  \widehat\pi_t(u_t;\alpha)
  =
  \frac{\phi_t(u_t;\alpha)^2+\tau_t(\alpha)\lambda_t(u_t)}
       {\widehat z_t(\alpha)},
\]
\[
  \widehat p_t^{\TT}(x_t,\vartheta;\alpha)
  =
  \int \widehat\pi_t(x_t,\vartheta,x_{t-1};\alpha)\,\dd x_{t-1}.
  \label{eq:p13-fixed-filter-output}
\]

\paragraph{Pseudocode.}
\begin{enumerate}[leftmargin=*]
\item Input the previous normalized TT filter
  \(\widehat p_{t-1}^{\TT}\), model densities \(f_\alpha,g_\alpha\), datum
  \(y_t\), and the frozen branch specification.
\item Build the target evaluation table.  For each frozen fitting point
  \(u^{(j)}=(x_t^{(j)},\vartheta^{(j)},x_{t-1}^{(j)})\), evaluate
  \[
    q_t(u^{(j)};\alpha)
    =
    \widehat p_{t-1}^{\TT}(x_{t-1}^{(j)},\vartheta^{(j)};\alpha)
    f_\alpha(x_t^{(j)}\mid x_{t-1}^{(j)},\vartheta^{(j)})
    g_\alpha(y_t\mid x_t^{(j)},\vartheta^{(j)}).
  \]
\item Apply the frozen coordinate maps and stabilizing shift.  Store
  \(s_t(u^{(j)};\alpha)=\exp(c_t/2)\sqrt{q_t(u^{(j)};\alpha)}\), or the
  equivalent shifted square-root target used by the branch.
\item Solve for TT cores core-by-core.  If the branch is interpolation-based,
  solve the frozen systems \(A_{t,k}g_{t,k}=b_{t,k}\).  If the branch is
  least-squares-based, solve the frozen normal equations
  \(A_{t,k}^{\top}W_{t,k}A_{t,k}g_{t,k}
  =A_{t,k}^{\top}W_{t,k}b_{t,k}\).  The branch must choose one of these value
  paths, or a precisely specified variant, before differentiation.
\item Assemble \(\phi_t(u)=G_{t,1}(u_1)\cdots G_{t,D}(u_D)\).  Contract the
  cores with the mass matrices to compute \(R_{t,0}\).
\item Set \(\widehat z_t=R_{t,0}+\tau_t\) and store
  \(\widehat\ell_t^{\TT}=\log\widehat z_t-c_t\).
\item Define the normalized joint approximation
  \[
    \widehat\pi_t(u)=
    \{\phi_t(u)^2+\tau_t\lambda_t(u)\}/\widehat z_t.
  \]
\item Marginalize the old state to obtain
  \[
    \widehat p_t^{\TT}(x_t,\vartheta)=
    \int\widehat\pi_t(x_t,\vartheta,x_{t-1})\,\dd x_{t-1}.
  \]
  Store the marginal numerator \(a_t(x_t,\vartheta)\) before division by
  \(\widehat z_t\), because the next derivative step needs
  \(\partial_i\log\widehat p_t^{\TT}\).
\item Store the cores, fitting matrices \(A_{t,k}\), right sides \(b_{t,k}\),
  weights \(W_{t,k}\) if used, target evaluations, coordinate maps, mass
  contractions \(R_{t,k}\), marginal contractions, numerator \(a_t\), shift
  \(c_t\), defensive mass \(\tau_t\), and any factorization signs or active
  ranks consumed by the derivative pass.
\item Output \(\widehat p_t^{\TT}\), \(\widehat\ell_t^{\TT}\), and the saved
  branch data.
\end{enumerate}

\section{Why We Need Two Propositions}

There are two different mathematical questions.

First, after replacing \(q_t\) by \(\phi_t^2+\tau_t\lambda_t\), do we still
have a legitimate filtering recursion?  A filter must be nonnegative, must
integrate to one, and must be carried forward in time by prediction, update,
normalization, and marginalization.  Proposition 1 answers exactly that
question for the fixed-branch approximation.  It does not say the approximation
is accurate; it says the approximation is a proper normalized filter.

Second, if we use
\[
  \widehat\ell_T^{\TT}(\alpha)
  =
  \sum_{t=1}^T\{\log\widehat z_t(\alpha)-c_t(\alpha)\}
\]
as a declared approximate likelihood, does the proposed gradient differentiate
that same scalar?  Proposition 2 answers that question.  It does not
differentiate the globally adaptive code; it differentiates the frozen branch
whose forward pass produced the scalar.

\section{Proposition 1: The Fixed Branch Still Filters}

\paragraph{Plain-language claim.}
If the previous approximate filter is a normalized density, and if the fixed
branch creates a nonnegative finite joint approximation
\(\widehat q_t=\phi_t^2+\tau_t\lambda_t\), then normalizing and marginalizing
it gives another normalized density.  Repeating this step gives a sequence of
proper approximate filters.

\paragraph{Objects being approximated.}
The exact object is \(q_t\) in \eqref{eq:p13-exact-q}.  The fixed branch uses
\[
  \widehat q_t(u_t;\alpha)
  =
  \phi_t(u_t;\alpha)^2+\tau_t(\alpha)\lambda_t(u_t)
\]
instead.  The filter is obtained from the normalized marginal of
\(\widehat q_t\).

\begin{assumption}[Dominated fixed-branch recursion]
For each \(t\), the previous approximate filter
\(\widehat p_{t-1}^{\TT}\) is nonnegative and integrates to one.  The
transition and observation densities are nonnegative and measurable.  The
fixed branch produces a measurable \(\phi_t\).  The defensive density satisfies
\(\lambda_t\ge0\) and \(\int\lambda_t=1\).  The defensive mass satisfies
\(\tau_t\ge0\).  Finally,
\[
  0<\widehat z_t(\alpha)
  =
  \int\{\phi_t(u_t;\alpha)^2+\tau_t(\alpha)\lambda_t(u_t)\}\,\dd u_t
  <\infty.
\]
\end{assumption}

\begin{proposition}[Fixed-branch squared-TT recursion defines a normalized approximate filter]
Define
\[
  \widehat\pi_t(u_t;\alpha)
  =
  \frac{\phi_t(u_t;\alpha)^2+\tau_t(\alpha)\lambda_t(u_t)}
       {\widehat z_t(\alpha)}
\]
and
\[
  \widehat p_t^{\TT}(x_t,\vartheta;\alpha)
  =
  \int
  \widehat\pi_t(x_t,\vartheta,x_{t-1};\alpha)\,\dd x_{t-1}.
\]
Under the dominated fixed-branch assumptions,
\(\widehat p_t^{\TT}\) is a nonnegative normalized density.  By induction from
a normalized prior, the recursion defines a sequence of normalized
approximate filtering densities.
\end{proposition}

\begin{proof}
The square \(\phi_t^2\) is nonnegative.  The defensive term
\(\tau_t\lambda_t\) is also nonnegative because \(\tau_t\ge0\) and
\(\lambda_t\ge0\).  Therefore \(\widehat q_t=\phi_t^2+\tau_t\lambda_t\ge0\).
The normalizer \(\widehat z_t\) is finite and positive by assumption, so
\(\widehat\pi_t=\widehat q_t/\widehat z_t\) is well defined and nonnegative.
Also
\[
  \int\widehat\pi_t(u_t;\alpha)\,\dd u_t
  =
  \frac{1}{\widehat z_t(\alpha)}
  \int \widehat q_t(u_t;\alpha)\,\dd u_t
  =
  1.
\]
Now integrate the marginal:
\[
\begin{aligned}
  \int\widehat p_t^{\TT}(x_t,\vartheta;\alpha)
  \,\dd x_t\,\dd\vartheta
  &=
  \int\int
  \widehat\pi_t(x_t,\vartheta,x_{t-1};\alpha)
  \,\dd x_{t-1}\,\dd x_t\,\dd\vartheta\\
  &=
  \int\widehat\pi_t(u_t;\alpha)\,\dd u_t
  =
  1.
\end{aligned}
\]
Tonelli's theorem justifies exchanging the order of integration because the
integrand is nonnegative.  Hence the new marginal is nonnegative and
normalized.  The prior gives the base case, and the same argument advances the
filter from \(t-1\) to \(t\).  Induction proves the sequence claim.
\end{proof}

\paragraph{What this proposition gives us.}
It proves that the fixed-branch approximation is still a valid filtering
recursion in the basic probabilistic sense: each output is a nonnegative
normalized density.

\paragraph{What this proposition does not give us.}
It does not prove that \(\widehat p_t^{\TT}\) is close to the exact filter.
Exact recovery occurs only in the special case
\(\widehat q_t=q_t\) almost everywhere at every time step.

\section{Proposition 2: The Gradient Differentiates The Same Scalar}

\paragraph{Plain-language claim.}
If the branch is fixed and all the displayed computations are differentiable,
then differentiating the TT cores, mass contractions, defensive mass, and
stabilizing constant gives the exact derivative of the scalar that the forward
pass computed:
\[
  \widehat\ell_T^{\TT}(\alpha)
  =
  \sum_{t=1}^T\{\log\widehat z_t(\alpha)-c_t(\alpha)\}.
\]

\paragraph{Implementation path.}
The forward pass computes \(\widehat z_t=R_{t,0}+\tau_t\), where \(R_{t,0}\)
is the mass contraction of \(\phi_t^2\).  The derivative pass must therefore
differentiate the same contraction \(R_{t,0}\), not a different approximation.
It also must differentiate the previous filter because the next target
\(q_t=\widehat p_{t-1}^{\TT}fg\) depends on that previous approximate filter.
This is the practical meaning of ``same scalar'': the backward pass must
differentiate the stored forward branch that produced
\(\widehat\ell_T^{\TT}\).  If the derivative is compared against a fresh
adaptive rerun with different ranks, pivots, samples, or factorization signs,
the comparison is no longer testing the derivative of the same formula.

\begin{assumption}[Fixed-branch differentiability]
The assumptions of Proposition 1 hold.  The branch choices are fixed.  Along
that branch, the model densities, previous approximate filter, coordinate
maps, basis maps, TT core coefficients, defensive mass \(\tau_t\), and
stabilizing constant \(c_t\) are differentiable in the parameter coordinate
\(\alpha_i\).  Differentiation may be interchanged with the finite contractions
and the displayed integrals.  Fixed interpolation or least-squares systems are
nonsingular, or a specific differentiable pseudoinverse branch is declared.
\end{assumption}

\begin{proposition}[Analytical gradient of the declared fixed-branch scalar]
Under the fixed-branch differentiability assumptions,
\[
  \partial_i\widehat\ell_T^{\TT}(\alpha)
  =
  \sum_{t=1}^T
  \left[
  \frac{\partial_i\widehat z_t(\alpha)}{\widehat z_t(\alpha)}
  -
  \partial_i c_t(\alpha)
  \right],
  \label{eq:p13-score-general}
\]
where, for fixed normalized \(\lambda_t\),
\[
  \partial_i\widehat z_t
  =
  2\int
  \phi_t(u;\alpha)\partial_i\phi_t(u;\alpha)\,\dd u
  +
  \partial_i\tau_t.
  \label{eq:p13-zdot-integral}
\]
If \(\phi_t\) is represented by TT cores and \(R_{t,0}\) is computed by the
mass contraction \eqref{eq:p13-mass-contraction}, then the direct
mass-contraction score is
\[
  \boxed{
  \partial_i\widehat\ell_T^{\TT}
  =
  \sum_{t=1}^T
  \left[
  \frac{\partial_iR_{t,0}+\partial_i\tau_t}
       {R_{t,0}+\tau_t}
  -
  \partial_i c_t
  \right].
  }
  \label{eq:p13-final-score}
\]
\end{proposition}

\begin{proof}
Because the branch is fixed, the scalar is a finite sum of differentiable
terms.  For each \(t\), \(\widehat z_t>0\), so
\[
  \partial_i\{\log\widehat z_t-c_t\}
  =
  \frac{\partial_i\widehat z_t}{\widehat z_t}-\partial_i c_t.
\]
Summing over \(t\) gives \eqref{eq:p13-score-general}.

Next,
\[
  \widehat z_t
  =
  \int\{\phi_t(u;\alpha)^2+\tau_t(\alpha)\lambda_t(u)\}\,\dd u.
\]
The assumed interchange of differentiation and integration gives
\[
  \partial_i\widehat z_t
  =
  \int 2\phi_t(u;\alpha)\partial_i\phi_t(u;\alpha)\,\dd u
  +
  \partial_i\tau_t\int\lambda_t(u)\,\dd u.
\]
Since \(\int\lambda_t=1\), this proves \eqref{eq:p13-zdot-integral}.  If
\(\lambda_t\) itself depends on \(\alpha_i\), an extra term
\(\tau_t\int\partial_i\lambda_t\,\dd u\) appears; it vanishes only when the
defensive density remains normalized on the differentiated branch.

It remains to connect \(\partial_i\phi_t\) to the stored TT computation.  If
\[
  \phi_t(u)=G_{t,1}(u_1)\cdots G_{t,D}(u_D),
\]
then the product rule gives
\[
  \partial_i\phi_t(u)
  =
  \sum_{k=1}^D
  G_{t,1}(u_1)\cdots
  \partial_iG_{t,k}(u_k)\cdots
  G_{t,D}(u_D).
  \label{eq:p13-tt-product-rule}
\]
Using the basis expansion
\[
  G_{t,k}(u_k)=\sum_{a=1}^{n_k}C_{t,k,a}b_{k,a}(u_k),
\]
the derivative \(\partial_iG_{t,k}\) is obtained by replacing
\(C_{t,k,a}\) with \(\partial_iC_{t,k,a}\), plus any basis-map derivative if
the coordinate map or basis depends smoothly on \(\alpha_i\).

The derivative of the mass contraction follows by differentiating the same
backward recursion that computed \(R_{t,0}\).  Starting from
\(\partial_iR_{t,D}=0\),
\[
\begin{aligned}
  \partial_i R_{t,k-1}
  =
  \sum_{a,b=1}^{n_k}M_k[a,b]\Big[
  &(\partial_iC_{t,k,a})R_{t,k}C_{t,k,b}^\top
  +C_{t,k,a}(\partial_iR_{t,k})C_{t,k,b}^\top\\
  &+C_{t,k,a}R_{t,k}(\partial_iC_{t,k,b})^\top
  \Big].
  \label{eq:p13-rdot-contraction}
\end{aligned}
\]
This is the matrix product rule applied to the contraction formula
\eqref{eq:p13-mass-contraction}.  Therefore
\(\partial_iR_{t,0}\) is the derivative of the square mass used in the forward
normalizer.

The core sensitivities come from the fixed core-construction equations.  If a
core vector \(g_{t,k}\) is determined by a frozen interpolation equation
\[
  A_{t,k}(\alpha)g_{t,k}(\alpha)=b_{t,k}(\alpha),
\]
then
\[
  A_{t,k}\partial_i g_{t,k}
  =
  \partial_i b_{t,k}
  -
  (\partial_i A_{t,k})g_{t,k}.
  \label{eq:p13-interp-sensitivity}
\]
If a core vector is determined by a fixed weighted least-squares problem
\[
  g_{t,k}
  =
  \arg\min_g
  \|W_{t,k}^{1/2}(A_{t,k}g-b_{t,k})\|_2^2
\]
with fixed \(W_{t,k}\) and nonsingular
\(N_{t,k}=A_{t,k}^\top W_{t,k}A_{t,k}\), differentiating the normal equations
gives
\[
\begin{aligned}
  N_{t,k}\partial_i g_{t,k}
  =
  &(\partial_i A_{t,k})^\top W_{t,k}
    (b_{t,k}-A_{t,k}g_{t,k})\\
  &+
  A_{t,k}^\top W_{t,k}
  (\partial_i b_{t,k}-(\partial_iA_{t,k})g_{t,k}).
  \label{eq:p13-ls-sensitivity}
\end{aligned}
\]
Thus the derivative pass solves the differentiated version of the same
linear equations used by the forward pass.

Finally, the target values \(b_{t,k}\) depend on the previous filter.  From
\[
  q_t
  =
  \widehat p_{t-1}^{\TT}f_\alpha g_\alpha,
\]
one has, wherever the factors are positive,
\[
  \partial_i\log q_t
  =
  \partial_i\log\widehat p_{t-1}^{\TT}
  +
  \partial_i\log f_\alpha
  +
  \partial_i\log g_\alpha.
  \label{eq:p13-logq-sensitivity}
\]
If the previous filter can be written as
\[
  \widehat p_{t-1}^{\TT}(v;\alpha)
  =
  \frac{a_{t-1}(v;\alpha)}{\widehat z_{t-1}(\alpha)},
\]
then
\[
  \partial_i\log\widehat p_{t-1}^{\TT}(v;\alpha)
  =
  \frac{\partial_i a_{t-1}(v;\alpha)}{a_{t-1}(v;\alpha)}
  -
  \frac{\partial_i\widehat z_{t-1}(\alpha)}
       {\widehat z_{t-1}(\alpha)}.
  \label{eq:p13-prev-filter-sensitivity}
\]
The numerator \(a_{t-1}\) is a marginal contraction of the previous
squared-TT object.  Its derivative is obtained by the same marginal
contraction with one previous core at a time replaced by its differentiated
core, plus the derivative of \(\tau_{t-1}\) and any defensive-density term.
Therefore the sensitivity recursion follows the same time recursion as the
filter itself.

Combining the log-normalizer derivative with the differentiated contraction
gives the boxed score \eqref{eq:p13-final-score}.  Every term in that formula
comes from differentiating the forward scalar
\(\sum_t\{\log\widehat z_t-c_t\}\), so the gradient is a same-scalar gradient
for the fixed branch.
\end{proof}

\paragraph{What this proposition gives us.}
It gives an implementable analytical derivative for the declared fixed-branch
approximate likelihood.  A prototype can compute the forward scalar, save the
branch data, compute core sensitivities, propagate mass-contraction
sensitivities, propagate previous-filter sensitivities, and assemble
\eqref{eq:p13-final-score}.

\paragraph{What this proposition does not give us.}
It does not give a global derivative of an adaptive algorithm whose ranks,
pivots, truncations, sample sets, or factorization branches change with
\(\alpha\).  It also does not prove that this approximate likelihood is close
to the exact likelihood.

\section{How To Compute The Gradient}

The following recipe is the minimal implementation contract for the
fixed-branch gradient.

\paragraph{Inputs.}
The implementation needs the state-space model, data, prior, frozen branch
specification, defensive densities, defensive masses, basis and mass matrices,
and derivatives of \(f_\alpha\), \(g_\alpha\), coordinate maps, shifts
\(c_t\), and any parameter-dependent branch-smooth objects.

\paragraph{Forward pass.}
For \(t=1,\ldots,T\):
\begin{enumerate}[leftmargin=*]
\item build the target evaluations of
  \(q_t=\widehat p_{t-1}^{\TT}f_\alpha g_\alpha\);
\item fit the square-root TT cores using the frozen equations;
\item compute \(R_{t,0}\), \(\widehat z_t=R_{t,0}+\tau_t\), and
  \(\widehat\ell_t^{\TT}=\log\widehat z_t-c_t\);
\item normalize and marginalize to obtain \(\widehat p_t^{\TT}\);
\item save all contractions and solve matrices needed by the derivative pass.
\end{enumerate}

\paragraph{Sensitivity pass.}
For each parameter coordinate \(\alpha_i\):
\begin{enumerate}[leftmargin=*]
\item differentiate the target evaluations using the displayed log-target
  identity
  \(\partial_i\log q_t
  =
  \partial_i\log\widehat p_{t-1}^{\TT}
  +\partial_i\log f_\alpha+\partial_i\log g_\alpha\).
\item solve the displayed interpolation sensitivity equation
  \(A_{t,k}\partial_i g_{t,k}
  =
  \partial_i b_{t,k}-(\partial_iA_{t,k})g_{t,k}\), or the displayed weighted
  least-squares sensitivity equation based on
  \(N_{t,k}=A_{t,k}^{\top}W_{t,k}A_{t,k}\).
\item compute \(\partial_iR_{t,0}\) by differentiating the same backward
  mass-contraction recursion used to compute \(R_{t,0}\), replacing one core
  factor at a time by its differentiated core and propagating
  \(\partial_iR_{t,k}\) backward.
\item compute \(\partial_i\widehat z_t=\partial_iR_{t,0}+\partial_i\tau_t\);
\item compute the step score
  \[
    \partial_i\widehat\ell_t^{\TT}
    =
    \frac{\partial_i\widehat z_t}{\widehat z_t}
    -
    \partial_i c_t;
  \]
\item differentiate the normalized marginal filter using
  \[
    \partial_i(\widehat q_t/\widehat z_t)
    =
    \frac{\partial_i\widehat q_t}{\widehat z_t}
    -
    \frac{\widehat q_t\,\partial_i\widehat z_t}{\widehat z_t^2},
  \]
  then marginalize to carry \(\partial_i\widehat p_t^{\TT}\) forward.
\end{enumerate}

\paragraph{Parity test before using the gradient.}
For each coordinate, freeze the same branch and compare
\[
  \frac{
  \widehat\ell_T^{\TT}(\alpha+h e_i)
  -
  \widehat\ell_T^{\TT}(\alpha-h e_i)}
  {2h}
  \quad\text{with}\quad
  \partial_i\widehat\ell_T^{\TT}(\alpha).
\]
If the finite-difference calls change a rank, pivot, sample set, sign,
truncation, or selected shift rule, the test has detected branch instability.
That is not a proof that TT filtering fails; it is a warning that the frozen
branch gradient no longer matches the adaptive value path being tested.

\section{What This Construction Gives Us}

The construction gives a clear approximate filtering object.  At each time,
it represents the unnormalized filtering joint density by a nonnegative
squared-TT density, computes its normalizer, and marginalizes the old state.
This is enough to produce a sequence of normalized approximate filters.

It also gives a declared approximate likelihood:
\[
  \widehat\ell_T^{\TT}
  =
  \sum_t\{\log\widehat z_t-c_t\}.
\]
For a fixed branch, the gradient in Proposition 2 is the derivative of that
same scalar.  This is the mathematical prerequisite for treating the method as
a future gradient-based likelihood candidate: one knows which scalar is being
differentiated and which operations must be differentiated.

The construction also gives a practical bridge from the published adaptive
method to an implementable prototype.  One can first implement the frozen
branch, check value-gradient parity, and then study how often the adaptive
branch changes in realistic filtering problems.

\section{What Still Must Be Tested}

Several important questions remain numerical rather than algebraic.

\begin{itemize}[leftmargin=*]
\item Approximation accuracy: the propositions do not say that
  \(\phi_t^2+\tau_t\lambda_t\) is close to \(q_t\).
\item Branch stability: finite-difference perturbations may change ranks,
  pivots, sample points, truncations, or factorization signs.
\item Conditioning: interpolation and least-squares systems can be ill
  conditioned, making both values and derivatives unstable.
\item Scaling: the method is persuasive for high dimension only if the TT ranks
  remain moderate on the target model class.
\item Comparison with fixed sparse-grid filtering: TT filtering is attractive
  when density or transport structure has stable low TT rank; sparse-grid
  filtering is attractive when the effective quadrature dimension is small.
\end{itemize}
These are not reasons to discard the method.  They are the experiments and
diagnostics needed before claiming that this approximation is reliable on a
specific high-dimensional nonlinear state-space model.

\appendix
\section{Source And Code Anchors}

The main text derives the filtering objects before citing them.  The source
anchors below record where the same technical objects appear in the inspected
literature and audit snapshot.

\begin{itemize}[leftmargin=*]
\item Zhao--Cui equations (1)--(3) define the state-space transition,
  observation, and parameter notation used by their sequential learning
  problem \cite{zhao-cui-2024}.
\item Zhao--Cui equations (9)--(12) and Algorithm 1 form the nonseparable
  sequential density over \((x_t,\theta,x_{t-1})\), approximate it in TT form,
  and integrate it.
\item Zhao--Cui equation (13) and nearby Lemma 1 introduce the squared-TT
  defensive density of the form \(\phi^2+\tau\lambda\).
\item Zhao--Cui Algorithm 2 applies the squared-TT construction sequentially.
\item Zhao--Cui Section 4.1 discusses the exact joint normalizer as evidence.
\item Zhao--Cui Propositions 2 and 4 support the marginal and conditional
  constructions used for filtering and conditional KR maps.
\item Cui--Dolgov Sections 2--3 provide the inverse Rosenblatt, functional TT,
  and squared inverse Rosenblatt transport substrate.
\item The audit snapshot of the companion code shows the same scalar path.
  In \texttt{models/full\_sol.m}, the running marginal likelihood receives
  the increment \texttt{log(sirt.z)-const}.  In the TT marginalization routine,
  the stored normalizer is \texttt{obj.z = obj.fun\_z + obj.tau}.  These code
  references show that the inspected MATLAB snapshot computes the same scalar
  structure; the mathematical justification remains the derivations in the
  main text.
\end{itemize}

\section{A Compact Checklist For A Prototype}

\begin{longtable}{@{}p{0.28\linewidth}p{0.62\linewidth}@{}}
\toprule
Item & What must be specified \\
\midrule
\endhead
Model & Prior, transition, observation density, data, parameter coordinate,
and all dimensions. \\
Branch & Coordinate order, domains, bases, ranks, fitting points, solver
schedule, defensive density, defensive mass, and shift rule. \\
Forward scalar & \(R_{t,0}\), \(\widehat z_t=R_{t,0}+\tau_t\), and
\(\widehat\ell_t^{\TT}=\log\widehat z_t-c_t\). \\
Filter output & Normalized joint density and marginal
\(\widehat p_t^{\TT}(x_t,\vartheta)\). \\
Derivative & Target sensitivities, core sensitivities, contraction
sensitivities, previous-filter sensitivities, \(\partial_i\tau_t\), and
\(\partial_i c_t\). \\
Diagnostics & Positive finite normalizers, stable branch objects,
conditioning, rank growth, and same-scalar finite-difference parity. \\
\bottomrule
\end{longtable}

\begin{thebibliography}{9}
\bibitem{zhao-cui-2024}
Zhao, Y. and Cui, T. (2024).
\newblock Tensor-train methods for sequential state and parameter learning in
state-space models.
\newblock \emph{Journal of Machine Learning Research}, 25(244):1--51.

\bibitem{cui-dolgov-2022}
Cui, T. and Dolgov, S. (2022).
\newblock Deep composition of tensor-trains using squared inverse Rosenblatt
transports.
\newblock \emph{Foundations of Computational Mathematics}, 22:1863--1922.
\newblock DOI: 10.1007/s10208-021-09537-5.
\end{thebibliography}

\end{document}
